\documentclass{article}
\usepackage[utf8]{inputenc}
\usepackage[spanish, es-tabla]{babel} % es-tabla para "Tabla" en lugar de "Cuadro"
\usepackage{amsmath}
\usepackage{geometry}
\usepackage{booktabs}
\usepackage{siunitx}
\usepackage{circuitikz}
\usepackage{enumitem} % Mejor que enumerate
\usepackage{caption}
\usepackage{makecell} % Para encabezados de tabla

\geometry{a4paper, margin=1in}

\begin{document}

\noindent\hfill\textbf{146 \hspace{1cm} Prácticas de electricidad}\par
\vspace{0.5cm}

    % Figura 
    \begin{center}
        \begin{circuitikz}[scale=0.9, transform shape]
            \draw (0,3) to[V, l_=$10\,\text{V}$, v=$V$] (0,0);
            \draw (0,3) to[R, l=$R$, l_=$1\,\text{k}\Omega$] (3,3);
            \draw (3,3)
                  to[short, -o] (4,3) node[right] {SIB};
            \draw (4,3) to[short, o-*] (4,2);
            \node[right] at (4.2, 2) {A};
            
            \draw (0,0) to[short] (3,0);
            \draw (3,0)
                  to[short, -o] (4,0) node[right] {SIA};
            \draw (4,0) to[short, o-*] (4,1);
            \node[right] at (4.2, 1) {B};
    
            \draw (4,2) to[vR, l=$R_L$, l_=$10\,\text{k}\Omega$] (4,1);
    
            \draw[dashed] (3.5,3.2) -- (3.5,-0.2);
        \end{circuitikz}
        \captionof{figure}{Circuito experimental para determinar la máxima transferencia de potencia a una carga.}
        \label{fig:circuito-24-1}
    \end{center}

    \begin{enumerate}[label=\arabic*., start=1]
        \item ...mente acoplados). Ajustar la tensión $V = 10$ voltios y mantenerla en este nivel durante el resto de la práctica.
        \item Abrir el interruptor $S_1$. Ajustar el reóstato para que $R_L$ sea 0 ohmios, medidos con un óhmetro. Cerrar $S_1$. Medir con un voltímetro electrónico la tensión $V_L$ entre los extremos de $R_L$ y anotar el resultado en la tabla 24-2.
        \item Abrir el interruptor $S_1$. Ajustar el reóstato de modo que $R_L$ mida 100 ohmios. Cerrar $S_1$. Medir con un EVM la tensión $V_L$ entre los extremos de $R_L$ y anotarla en la tabla 24-2.
        \item Repetir la operación 3 para cada valor de $R_L$ indicado en la tabla.
        \item Por los valores medidos correspondientes de $V_L$ y $R_L$, calcular la potencia de $R_L$ utilizando la fórmula $W = V_L^2/R_L$. Convertir W en milivatios y anotar en la tabla 24-2. Presentar una muestra del cálculo.
        \item Calcular $W_T$ utilizando la fórmula $W_T = V^2/(R + R_L)^2$. Convertir $W_T$ en milivatios y anotar en la tabla 24-2. Presentar una muestra del cálculo.
        \item Dibujar un gráfico de W en función de $R_L$, correspondiendo $R_L$ al eje de abscisas u horizontal, y W al de ordenadas o vertical.
    \end{enumerate}
    

\begin{center}
    \textbf{TABLA 24-2}
\end{center}

\sisetup{output-decimal-marker={,}}
\centering
\begin{tabular}{
    S[table-format=5.0]
    S[table-format=5.0]
    S[table-format=1.3]
    S[table-format=2.3]
    S[table-format=2.3]
}
\toprule
{\makecell{\textbf{$R_L$,} \\ \textbf{ohmios}}} &
{\makecell{\textbf{$R+R_L$,} \\ \textbf{ohmios}}} &
{\makecell{\textbf{$V_L$,} \\ \textbf{voltios}}} &
{\makecell{\textbf{$W = \frac{V_L^2}{R_L}$} \\ \textbf{(milivatios)}}} &
{\makecell{\textbf{$W_T = \frac{V^2}{(R+R_L)^2}$} \\ \textbf{(milivatios)}}} \\
\midrule
0 & & & & \\
100 & & & & \\
200 & & & & \\
400 & & & & \\
600 & & & & \\
800 & & & & \\
850 & & & & \\
900 & & & & \\
950 & & & & \\
1000 & & & & \\
1100 & & & & \\
1200 & & & & \\
1500 & & & & \\
1700 & & & & \\
2000 & & & & \\
4000 & & & & \\
6000 & & & & \\
8000 & & & & \\
10000 & & & & \\
\bottomrule
\end{tabular}

\end{document}